\section{Problem Statement: STOP as a Diagnostic Outcome}

In enterprise practice, a decision to halt, suspend, or refuse further progression of an initiative is commonly interpreted as a failure. This interpretation is usually implicit rather than argued: STOP is taken to indicate insufficient analysis, lack of data, weak stakeholder alignment, or analyst incapacity. Within this framing, continuation is treated as the default success state of analysis, while STOP is relegated to an exceptional or undesirable outcome. Such assumptions align with classical views of decision-making that presuppose the availability of improved choices given sufficient analysis, despite well-known limits of rationality in complex organizational contexts \cite{simon1957}.

This paper questions the validity of that assumption from a formal modeling perspective. When enterprises are treated as structured systems evolving under constraints, the outcome of an analysis cannot be evaluated solely by whether it recommends continuation. The central issue is not whether an analysis produces forward motion, but whether the modeled system admits a continuation that is internally consistent with its constraints, structure, and admissible state space. This systems-oriented perspective follows the tradition of general system theory, in which global behavior cannot be reduced to isolated local decisions \cite{bertalanffy1968}.

The diagnostic paradox addressed here can be stated as follows. In operational reality, analysts are frequently confronted with situations in which all locally available actions appear feasible in isolation, yet are globally incompatible when considered together. In such cases, forcing a recommendation to proceed may preserve the appearance of decisiveness while violating the internal coherence of the system model. This tension between local feasibility and global incompatibility is a characteristic feature of complex systems, where aggregate behavior exhibits properties not inferable from individual components \cite{anderson1972}. Conversely, a STOP or No-Go outcome may reflect not analytical failure, but the absence of any admissible continuation under the current modeling assumptions.

From a formal standpoint, this shifts the locus of evaluation. STOP is not primarily a statement about analyst competence, effort, or data volume. It is a statement about model admissibility. If, under a given model of the enterprise, no continuation satisfies the model’s structural constraints, then continuation is not a valid analytical outcome, regardless of external pressure or organizational preference. In this sense, STOP functions as a diagnostic signal indicating that the current model does not support forward evolution, analogous to constraint-induced limits in cybernetic systems \cite{ashby1956}.

This paper therefore treats STOP as a legitimate analytical outcome arising from the interaction between enterprise structure and modeled constraints. The question is not whether STOP is desirable, but whether it is formally justified. The widespread equation of STOP with failure presupposes that admissible continuation always exists if analysis is performed competently. That presupposition itself is the object of examination.

The single frozen question addressed by this paper (Q2) is the following: under what formal conditions does a STOP or No-Go result represent a valid diagnostic outcome of an enterprise model, rather than a failure of analysis?

To address this question, the paper analyzes STOP within the framework of enterprises modeled as continua (K\_EA). It examines how structural constraints, admissibility of states, and internal coherence delimit the space of valid analytical outcomes. The focus is on distinguishing cases in which STOP emerges as a necessary consequence of the model from cases in which it would indeed indicate analytical insufficiency.

The paper will not provide prescriptions for governance, decision-making, or organizational behavior. It will not propose operational remedies, escalation paths, or managerial actions. It will not redefine STOP in practical or normative terms. Its scope is limited to establishing, at a formal level, whether and when STOP can be considered a valid diagnostic result within a continuum-based enterprise model.

By reframing STOP as a question of model admissibility rather than analytical failure, the paper aims to clarify the logical status of No-Go outcomes in enterprise analysis. This clarification is a prerequisite for any subsequent methodological or practical discussion, which is explicitly outside the scope of the present work.
